\documentclass[12pt]{article}

\usepackage[left=1.5in,right=1in,top=1in,bottom=1in]{geometry}
\usepackage[T1]{fontenc}
\usepackage[utf8]{inputenc}
\usepackage{lmodern}
\usepackage{microtype}
\usepackage{setspace}
\setstretch{1.15}
\setlength{\parindent}{0pt}
\setlength{\parskip}{0.75em}

\usepackage{titlesec}
\titleformat{\section}
  {\normalfont\Large\bfseries}
  {}
  {0pt}
  {}
\titlespacing*{\section}{0pt}{3ex}{1.5ex}

\usepackage{fancyhdr}
\pagestyle{fancy}
\fancyhf{}
\fancyhead[L]{\textit{Flower Wars -- Interpretive Essay}}
\fancyhead[R]{\thepage}

\title{\textbf{Bounded Violence}\\
\large An Interpretive Essay on \textit{Flower Wars}}
\author{Flyxion}
\date{\today}

\begin{document}
\maketitle

\section*{Introduction}

\textit{Flower Wars} is not, at its core, a historical drama about the Mexica empire. It is an inquiry into whether violence can be made locally tractable without becoming globally catastrophic. The historical setting provides a rare and unusually explicit case study of an attempted solution: the ritualization, calendrical scheduling, and mutual recognition of war as a bounded practice. This essay interprets the screenplay through a set of theoretical commitments that recur across my broader work: constraint-before-content, local solvability versus global dysfunction, and the idea that institutions are not abstractions but embodied control systems whose failure modes are structural rather than moral.

The Flower Wars function in the narrative as a deliberately engineered constraint surface. They are not peace. They are an attempt to transform unconstrained gradient descent into a controlled oscillation. Violence is not eliminated; it is phase-locked.

\section*{Provenance and First Encounter}

The origin of this project is not academic, and not clean. It begins with a book that was already decaying.

I was given a moldy paperback by an old man on an island off the coast of Mexico. It was not offered ceremoniously. It emerged from a stack of other books--sun-warped, salt-stained, pages swelling slightly from years of humidity. The cover was faded, the glue in the spine failing. It smelled unmistakably of mildew and time. He told me, casually, that it was about Tlacaelel. That was all.

The book was in Spanish. At the time, my Spanish was imperfect but serviceable, and the difficulty of reading it mattered. It slowed me down. It forced attention. The prose resisted skimming. I read it in fragments at first, then in longer stretches, often re-reading pages not because I was confused, but because the tone felt unfamiliar in a way that demanded reorientation.

The title was \textit{Tlacaelel: El azteca entre los aztecas}, by Antonio Velasco Piña, published in 1978. It was not a scholarly history. It was openly mythic, speculative, and devotional in places. It mixed historical figures with spiritual symbolism, political intrigue with mysticism, and it made no attempt to apologize for doing so. Tlacaelel was presented not merely as a strategist or reformer, but as a liminal figure--operating behind power, reshaping the empire without wearing its crown.

What held my attention was not the mysticism itself, but the structural implication beneath it. The novel repeatedly returned to the idea that Tlacaelel's power did not come from sovereignty, but from architecture: from shaping beliefs, rituals, and institutions that outlived individual rulers. Even where the author drifted into fantasy, the underlying intuition remained sharp. Power did not flow from command. It flowed from constraint.

Reading the book in that condition--physically deteriorating, linguistically resistant, temporally distant--produced a particular kind of immersion. It did not feel like consuming a story. It felt like encountering a residue. Something half-preserved, half-decayed, carrying an idea that had survived in compromised form.

This matters, because \textit{Flower Wars} is not an adaptation of Velasco Piña's novel. It rejects much of its mysticism, its nationalist romanticism, and its heroic framing. But it inherits something quieter and more durable: the intuition that Tlacaelel was not important because he won battles, but because he attempted to redesign the conditions under which violence occurred.

The screenplay takes that intuition and strips it of prophecy, destiny, and spiritual exceptionalism. What remains is an engineering problem. How does a civilization manage irreversible harm without erasing itself? How does power operate when it is exercised through fields, calendars, and memory rather than decree?

In retrospect, it feels appropriate that this project began with a book already coming apart. The story it told was itself an unstable mixture of history and invention, devotion and distortion. But inside that instability was a signal strong enough to survive translation, decay, and reinterpretation.

This film is my attempt to extract that signal, subject it to constraint-first thinking, and see what remains once myth is treated not as truth, but as an interface through which older civilizations grappled with irreversibility.

The mold did not matter. The island did not matter. Even the accuracy of the novel did not finally matter.

What mattered was the encounter with an idea that had not yet exhausted itself.

\section*{Constraint Before Content}

A central error in modern interpretations of both violence and governance is the assumption that meaning precedes structure. \textit{Flower Wars} inverts this assumption. Tlacaelel's intervention does not begin with ethical exhortation, theological revision, or appeals to mercy. It begins with accounting. Bodies, calendars, fertility rates, warrior replacement times, and cosmological obligations are treated as coupled variables in a single dynamical system.

This reflects a constraint-first worldview: outcomes are downstream of boundary conditions, not intentions. The Flower Wars do not rely on warriors becoming kinder. They rely on warriors being unable to act outside a defined field without consequence. The emphasis on markers, schedules, and enforcement mechanisms is therefore not incidental; it is the substance of the reform.

In my theoretical vocabulary, this is the difference between content-level alignment and constraint-level alignment. The former fails under pressure. The latter can persist even when individual actors defect, provided the cost of defection exceeds its advantage. The tragedy of the Flower Wars is not that warriors remain violent, but that the constraint regime cannot scale to actors who do not recognize its grammar.

\section*{Local Tractability and Global Breakdown}

The Flower Wars succeed precisely where most moral systems fail: locally. Casualties stabilize. Capture replaces annihilation. Memory replaces vendetta. At the scale of the field, the system is elegant. Its breakdown does not occur because it is incoherent, but because its coherence is local.

This distinction matters. Systems optimized for local tractability often appear stable until they encounter an interaction regime outside their design envelope. The arrival of the Spanish is not framed as evil overwhelming good, but as an incompatible control system entering the same space. The Spaniards do not violate the rules of the Flower Wars; they never enter them. There is no shared phase space.

This maps directly onto my broader concern with global dysfunction emerging from locally rational systems. The Mexica are not irrational. Tlacaelel is not naive. The failure occurs because bounded violence requires mutual recognition, and recognition is not enforceable against actors whose payoff structure lies elsewhere.

The film therefore resists counterfactual fantasies. There is no moment where a better argument, a faster reform, or a stronger enforcement mechanism would have solved the problem. The Flower Wars were not underpowered; they were mis-scaled.

\section*{Ritual as Control Surface}

Ritual in \textit{Flower Wars} is not superstition. It is an interface. Calendars, chants, sacrifices, and fields function as synchronization mechanisms across distributed agents. They align expectations, pace escalation, and provide shared reference points for meaning.

From this perspective, sacrifice is not portrayed as barbarism but as a crude yet effective signal amplifier. Blood is visible. Cost is undeniable. Ritual converts abstract obligation into embodied certainty. The High Priest Tizoc is therefore not an antagonist in the moral sense, but a rival systems engineer operating under different constraints. Where Tlacaelel seeks sustainability through bounded repetition, Tizoc seeks compliance through saturation.

The tension between them is not belief versus skepticism. It is between two control strategies: modulation versus shock. The escalation of ritual as the Flower Wars weaken mirrors a familiar failure mode in modern institutions, where increasing intensity is mistaken for increasing control.

\section*{Infrastructure, Memory, and Survivability}

One of the screenplay's quiet claims is that infrastructure outlives ideology. Runners continue after councils dissolve. Archives persist after policies fail. Fields disappear, but paths remain.

Xochitl's arc embodies this claim. She is not a theorist, priest, or ruler. She is a carrier. Her survival is not guaranteed by belief but by motion. When command structures collapse, she continues to run, reframing communication as an ethical act rather than an administrative one.

This reflects a deeper theoretical commitment: that survivable systems are those whose core functions degrade gracefully. The Flower Wars fail as policy but survive as memory. This survival is not heroic. It is archival. Teaching replaces enforcement. Description replaces command.

The final field is empty not because it was meaningless, but because meaning cannot be forced to persist. It must be re-entered voluntarily.

\section*{Why the Attempt Matters}

The temptation, especially in retrospect, is to dismiss the Flower Wars as a doomed experiment. This essay rejects that framing. Systems should not be evaluated solely by their permanence. Temporary success in constraining catastrophic behavior is itself a meaningful outcome.

In a universe where unbounded violence is the default attractor, even a partial, fragile, time-limited counterexample matters. The Flower Wars demonstrate that humans can construct institutions that interrupt escalation without requiring moral transformation. That this interruption is temporary does not negate its significance.

The deeper tragedy explored in \textit{Flower Wars} is not failure, but forgetting. The loss is not the system itself, but the memory that such a system once functioned at all.

\section*{Conclusion}

\textit{Flower Wars} should be read not as a lament for a fallen civilization, but as a warning about the limits of local solutions in a globally adversarial world. Constraint works. Ritual works. Bounded violence works. But none of these survive contact with actors who refuse the field.

The film's final gesture is therefore not despair but modesty. It refuses redemption arcs and permanent solutions. Instead, it offers a record of an attempt: careful, costly, imperfect, and briefly successful.

In my broader theoretical work, this is the only kind of success worth taking seriously.

\section*{Mapping to RSVP: Entropy, Constraint, and Field Dynamics}

Within the RSVP framework, the Flower Wars can be read as an attempt to engineer a metastable entropy gradient inside a violent system that would otherwise collapse into maximal dissipation. Unconstrained warfare corresponds to a rapid entropy descent: bodies, social trust, and institutional memory are consumed faster than they can be replenished. Tlacaelel's intervention introduces an artificial potential well. The ritual field, the calendar, and the rules of capture function analogously to scalar and vector constraints in RSVP: they redirect flow without eliminating it.

The Flower War field itself behaves like a localized plenum cell. Energy is permitted to circulate, but only along permitted trajectories. Violence becomes laminar rather than turbulent. Importantly, this is not equilibrium. RSVP rejects static balance in favor of managed gradients. The Flower Wars do not stop entropy production; they slow and shape it so that structure persists longer than it otherwise would.

From this perspective, sacrifice is not an anomaly but a pressure-release valve. It allows entropy to exit the system in a controlled fashion, preventing catastrophic rupture elsewhere. The failure occurs when a higher-energy external system enters the plenum without coupling to its constraint geometry. The Spanish do not merely overpower the system; they short-circuit it, injecting entropy along dimensions the field cannot represent.

\section*{Aspect Relegation and the Illusion of Moral Progress}

Aspect Relegation Theory provides a lens for understanding why the Flower Wars appear, to later observers, either barbaric or naive. Ritualized violence becomes invisible as violence once its parameters are relegated beneath conscious deliberation. Warriors trained within the system no longer experience capture as mercy but as normal operation. The ethical novelty of restraint is compressed into habit.

This relegation is not moral failure; it is the necessary mechanism by which any constraint-based system operates. When the Spanish arrive, they are perceived as monstrously violent precisely because their actions have not been relegated. Their violence remains high-resolution, uncompressed, and therefore cognitively shocking.

The film insists on a crucial asymmetry: what looks like moral regression is often a change in resolution, not intent. The Flower Wars fail not because they were morally thin, but because their ethical compression could not be re-expanded quickly enough when the operating regime changed.

\section*{Local Solvability as a Design Principle}

The screenplay deliberately foregrounds local solvability over universal correctness. The Flower Wars are never claimed to be optimal in a global sense. They are sufficient under known conditions. This aligns with a design philosophy that treats solvability as context-bound and provisional.

In RSVP terms, the system is tuned to a specific region of phase space. Within that region, it is robust. Outside it, the system exhibits brittle failure. The narrative rejects the idea that a more powerful constraint would have solved this problem. Increasing rigidity would only have increased fragility.

This is why Moctezuma's skepticism is structurally valid even when ethically troubling. He intuits that predictability lowers deterrence, and that deterrence is a form of negative entropy. The film does not refute him; it places his logic alongside Tlacaelel's and lets the incompatibility stand unresolved.

\section*{Infrastructure as Persistent Vector Field}

Xochitl's runner network maps cleanly onto RSVP's vector field intuition. Information flow precedes and outlasts institutional authority. Even when scalar constraints (fields, calendars, offices) dissolve, vector continuity remains. Motion persists where structure collapses.

This distinction explains why messengers, paths, and archives survive longer than palaces or councils. RSVP predicts that systems with low semantic load but high transport reliability degrade more gracefully. Xochitl's continued running is not symbolic defiance; it is structural inevitability. Communication is cheaper than control.

The death whistle's semantic shift--from command signal to mnemonic artifact--illustrates how vector fields can retain directionality even after their original potential function disappears. The sound no longer orders action; it preserves orientation.

\section*{Ritual as Phase-Locking Mechanism}

Ritual, in RSVP terms, is a phase-locking protocol. Calendars synchronize distributed agents. Repetition enforces coherence across time. The Flower Wars succeed insofar as they keep participants phase-aligned: everyone knows when, where, and how violence will occur.

When the Spanish arrive, phase coherence is lost. There is no shared oscillation to lock onto. This is why the Mexica response escalates ritual intensity rather than abandoning it. The system attempts to regain coherence by increasing signal amplitude, mistaking volume for alignment.

RSVP predicts this failure mode. When phase-locking breaks, systems often attempt to compensate through energy injection rather than structural adaptation. The Great Temple scenes are therefore not regressions into superstition but textbook control failure under external perturbation.

\section*{Memory as Residual Structure}

The final teaching scenes articulate a core RSVP claim: when scalar and vector fields collapse, topological memory can remain. The Flower Wars persist not as an institution but as a remembered configuration of possibility. This memory is low-energy, low-fidelity, but portable.

In RSVP language, this is a transition from dynamic constraint to static invariant. The field no longer functions, but its existence constrains future imagination. It proves that bounded violence once occupied a stable basin. That proof matters even if it cannot be reenacted wholesale.

This is why the final empty field is not nihilistic. It is a latent configuration space, awaiting new boundary conditions.

\section*{Constraint Without Guarantees}

Mapped through RSVP and related constraint-first theories, \textit{Flower Wars} becomes a study in the limits of engineered stability. It affirms that humans can shape destructive forces without abolishing them, but denies that such shaping can be universalized.

The film does not offer lessons in morality. It offers lessons in systems design under existential pressure. Constraint works. Ritual works. Memory works. None of them work forever.

What endures is not the solution, but the demonstration that solutions can exist.

That, within my theoretical framework, is sufficient justification for the attempt.

\section*{Event History, Irreversibility, and Architectural Memory}

The narrative logic of \textit{Flower Wars} aligns closely with event-historical and irreversibility-based architectures that recur across my broader work. These architectures reject reversible state transitions, equilibrium assumptions, and feed-based optimization in favor of explicit, irreversible event traces. The film can therefore be read as an allegory for an event-sourced civilization: one that records, constrains, and learns from irreversible actions rather than pretending they can be undone.

At its core, the Flower Wars transform violence from an unlogged state transition into an event with mandatory inscription. Capture replaces killing not because it is kinder, but because it preserves reversibility at the level of meaning while maintaining irreversibility at the level of action. A captured warrior cannot retroactively become uncaptured, but their survival allows future negotiation. This distinction mirrors event-historical design, where events are irreversible but system meaning is continuously recomputed atop the event log.

\section*{War as Event-Sourced Computation}

In an event-sourced architecture, the system's truth is not its current state but the ordered history of events that produced it. The Flower Wars instantiate this principle explicitly. Each engagement is calendared, witnessed, recorded, and archived. The archive scenes are not narrative decoration; they are the computational substrate of the system.

Unconstrained warfare corresponds to state-overwriting computation. Cities change hands, populations vanish, and the past is erased through annihilation. The Flower Wars replace this with append-only violence. Every encounter adds to history rather than deleting its participants. This preserves informational continuity even as conflict persists.

Tlacaelel's insistence on counting is therefore not bureaucratic pathology. It is an irreversible ledger strategy. Once blood is spilled on a sanctioned field, it becomes part of a shared record that constrains future action. Retaliation is bounded because the past cannot be denied without breaking the system entirely.

\section*{Irreversibility as Moral and Structural Fact}

A central error in modern political and technological systems is the fantasy of reversibility: the belief that damage can be rolled back, reputations reset, environments restored to prior baselines. \textit{Flower Wars} refuses this fantasy. Death is final. Trauma persists. Memory accumulates.

By converting killing into capture, the system does not make violence reversible; it shifts irreversibility from biological termination to social obligation. A captured warrior returns marked, obligated, changed. The event cannot be undone, but its consequences propagate differently.

This maps directly onto irreversibility architectures in which events are immutable and compensatory actions must occur downstream rather than through rollback. The Flower Wars are a civilization-scale compensation mechanism: harm is acknowledged, bounded, and rebalanced over time rather than denied.

\section*{Fields as Commit Boundaries}

The ritual field functions as a commit boundary in an irreversible system. Inside the field, actions are valid, logged, and enforceable. Outside it, violence becomes undefined behavior. This distinction explains why violations are treated as existential threats rather than tactical ones: they break the transaction model.

Sanction scenes in the screenplay mirror what happens when a distributed system encounters a non-idempotent write outside its transaction scope. The response is not moral outrage but escalation of constraint--secondary sanctions, narrower terms, stricter enforcement. The system hardens because it must preserve coherence of its history.

When the Spanish arrive, they do not merely attack; they operate entirely outside the commit model. Their actions cannot be reconciled with the event log. The system cannot compensate because it cannot even recognize the events as valid inputs. Collapse follows not from weakness, but from architectural mismatch.

\section*{Memory as Persistence Layer}

The archive, the codices, and the oral teaching scenes function as persistence layers. When the active system fails, the persistence layer remains. This is the precise reason Tlacaelel pivots from governance to teaching. He recognizes that the runtime is unstable, but the log can survive.

Xochitl's transformation from runner to memory-carrier reinforces this. She no longer transports commands; she transports continuity. Her continued motion maintains causal linkage between otherwise disconnected nodes.

This reflects an architectural principle: when control paths fail, data paths may still function. Memory, if preserved, allows reconstruction under new regimes. The film's insistence on writing ``ended'' rather than ``failed'' is an explicit design choice. It preserves interpretability of the log.

\section*{Irreversible Collapse and Graceful Degradation}

A defining feature of resilient irreversible systems is graceful degradation. Components fail without erasing the event history. The Flower Wars do not vanish instantly. They fragment. Fields are abandoned, but paths remain. Rules dissolve, but language persists.

The final scenes depict this degraded mode: no enforcement, no scheduling, no ritual authority--only memory and example. This is not a moral resolution but an architectural one. The system cannot be restored, but its invariant--the proof that bounded violence once functioned--remains accessible.

In my own work on irreversibility and event history, this is the desired failure mode. Systems should fail legibly. They should leave behind traces that prevent total epistemic collapse. \textit{Flower Wars} argues that civilizations, like software, should be designed with their own eventual failure in mind.

\section*{Against Rewind Culture}

The film implicitly critiques rewind culture: the belief that catastrophic events can be erased, sanitized, or rewritten. Tlacaelel's destruction of prior dynastic records is historically attested and morally ambiguous, but the screenplay reframes this act as a warning rather than an endorsement. Erasing history makes future constraint impossible.

The Flower Wars are an attempt to build a civilization that does not rely on forgetting in order to function. Their failure does not invalidate that attempt; it clarifies its necessity.

\section*{Conclusion: Architecture Over Redemption}

Read through the lens of event history and irreversibility, \textit{Flower Wars} becomes a meditation on architectural ethics. It asks not whether violence can be justified, but whether it can be logged, constrained, and remembered without destroying the substrate that records it.

The answer offered is provisional and limited, but nontrivial: yes, briefly, under specific conditions, with immense discipline. And when it fails, what matters is not restoration, but persistence of the trace.

This aligns with my broader position: civilizations do not need perfect solutions. They need architectures that remember what happened when solutions were tried. 

\section*{Further Synthesis: Constraint, Eventhood, and Civilizational Design}

This final synthesis tightens the correspondence between the narrative architecture of \textit{Flower Wars} and the broader theoretical program that motivates it: a unified view of civilization as an event-historical, constraint-first system operating under irreversible dynamics. What the screenplay dramatizes, and what the theory formalizes, is the same proposition viewed at different resolutions: that survival is not secured by correctness, virtue, or dominance, but by the disciplined management of irreversible transitions.

\section*{From Governance to Geometry}

Across my work, a recurring move is the replacement of institutional metaphors with geometric ones. Laws are not rules but boundaries. Ethics is not preference but curvature. Memory is not recollection but topology. \textit{Flower Wars} adopts this move implicitly. The ritual field is not a law enforced by authority; it is a shaped space that makes certain trajectories cheap and others expensive.

Tlacaelel does not govern by issuing commands. He reshapes the geometry in which commands operate. Warriors are not persuaded to restrain themselves; they are placed in a field where unrestrained action exits the domain of meaning altogether. This is the same move made in irreversibility architectures, where one does not prevent bad states by moralizing, but by removing rollback paths and making destructive transitions one-way and visible.

Seen this way, the Flower Wars are not a policy but a geometry of violence.

\section*{Eventhood as the Unit of Meaning}

Another throughline in my theoretical work is the insistence that events, not states, are the atomic units of meaning. States are compressions; events are facts. The screenplay enacts this distinction rigorously. A Flower War engagement is meaningful not because of its outcome, but because it occurred, was witnessed, and was recorded. The archive matters more than the battlefield.

This reframes victory entirely. Winning ceases to be a terminal state and becomes an intermediate annotation. What matters is not who stands last, but what can still be said afterward without contradiction. This is precisely the logic of event-sourced systems, where truth is reconstructed from history rather than asserted as a snapshot.

The collapse of the Flower Wars is therefore legible as an event failure rather than a state failure. The system does not arrive at a bad equilibrium; it encounters events it cannot represent. Meaning collapses before power does.

\section*{Irreversibility Without Apocalypse}

A subtle but important synthesis point concerns how irreversibility is treated. Many theories that foreground irreversibility drift toward apocalyptic thinking: once a boundary is crossed, everything is lost. \textit{Flower Wars} resists this. Irreversibility is everywhere, but it is not total.

Deaths are irreversible. Cultural collapse is irreversible. But memory, pedagogy, and example remain. This layered irreversibility mirrors my architectural stance that systems should be designed so that high-level meaning degrades more slowly than low-level structure. Buildings fall; paths remain. Institutions fail; stories persist.

This is why the film ends not with conquest or annihilation, but with teaching. Teaching is the slowest, cheapest, and most robust persistence mechanism available to civilizations that have lost enforcement capacity. It is not redemption. It is survivable residue.

\section*{Constraint as the Opposite of Optimization}

One of the most easily misunderstood aspects of the Flower Wars, both historically and narratively, is the assumption that they represent optimization: maximizing captives, minimizing deaths, stabilizing borders. The synthesis here is sharper. Constraint is not optimization. Constraint is the refusal to optimize beyond a survivable regime.

Optimization pushes systems toward edge cases. Constraint pulls them back toward interior regions of phase space where error is absorbable. Tlacaelel is not an optimizer. He is an anti-optimizer. He deliberately leaves value on the table--territory unconquered, enemies alive--because total extraction collapses the substrate that extraction depends on.

This resonates directly with my critiques of modern systems that confuse efficiency with intelligence. The Flower Wars are inefficient by design. That inefficiency is the feature that keeps the system legible.

\section*{Why the Failure Is Not a Refutation}

The final synthesis point addresses the most obvious objection: if the Flower Wars failed, what do they prove? The answer, both in the screenplay and in the theory, is that failure under external regime change is not evidence against internal coherence.

In irreversible systems, the correct question is not ``Does it last forever?'' but ``Does it create a stable basin under known conditions?'' The Flower Wars demonstrably did. Their failure mode--structural mismatch rather than internal contradiction--is exactly the failure mode one would predict for a locally coherent constraint system encountering an unaligned external actor.

This matters because it rescues the attempt from trivial dismissal. The experiment did not fail because it was naïve. It failed because it was specific. And specificity is the cost of tractability.

\section*{Civilizations as Event-Handling Systems}

Pulling all strands together, \textit{Flower Wars} can be read as a proposal--rendered in narrative form--for how civilizations should be evaluated: not by their moral narratives or ideological coherence, but by their event-handling capacity.

Can the system absorb violence without erasing itself?
Can it record harm without amplifying it?
Can it fail without deleting its own memory?

The Mexica system, as reimagined here, answers ``yes'' to all three--for a time. That time matters.

\section*{Closing Synthesis}

The deeper unity between the screenplay and my theoretical work lies here: both reject salvation narratives. There is no final alignment, no permanent constraint, no stable endpoint. There are only better and worse ways of traversing irreversible terrain.

\textit{Flower Wars} is not an argument for ritualized violence. It is an argument for architectures that acknowledge violence as an event, constrain it spatially and temporally, and refuse the lie of reversibility.

In that sense, the film is not about the past. It is about design principles for any system--political, technological, or civilizational--that must operate under conditions where mistakes cannot be undone and forgetting is fatal.

That synthesis, rather than the historical specifics, is the work's enduring claim.

\newpage
\begin{thebibliography}{99}

\bibitem{velascopina1978}
Velasco Piña, Antonio.
\textit{Tlacaelel: El azteca entre los aztecas}.
México: Editorial Diana, 1978.

\bibitem{leonportilla1992}
León-Portilla, Miguel.
\textit{Aztec Thought and Culture: A Study of the Ancient Nahuatl Mind}.
Translated by Jack Emory Davis.
Norman: University of Oklahoma Press, 1992.

\bibitem{leonportilla1969}
León-Portilla, Miguel.
\textit{La filosofía náhuatl estudiada en sus fuentes}.
México: Universidad Nacional Autónoma de México, 1969.

\bibitem{sahagun1950}
Sahagún, Bernardino de.
\textit{Florentine Codex: General History of the Things of New Spain}.
Translated by Arthur J. O. Anderson and Charles E. Dibble.
Santa Fe: School of American Research, 1950--1982.

\bibitem{durán1994}
Durán, Diego.
\textit{The History of the Indies of New Spain}.
Translated by Doris Heyden.
Norman: University of Oklahoma Press, 1994.

\bibitem{clendinnen1991}
Clendinnen, Inga.
\textit{Aztecs: An Interpretation}.
Cambridge: Cambridge University Press, 1991.

\bibitem{smith2012}
Smith, Michael E.
\textit{The Aztecs}.
3rd ed.
Oxford: Wiley-Blackwell, 2012.

\bibitem{hassig1988}
Hassig, Ross.
\textit{Aztec Warfare: Imperial Expansion and Political Control}.
Norman: University of Oklahoma Press, 1988.

\bibitem{hassig1992}
Hassig, Ross.
\textit{War and Society in Ancient Mesoamerica}.
Berkeley: University of California Press, 1992.

\bibitem{graulich1997}
Graulich, Michel.
\textit{Myths of Ancient Mexico}.
Translated by Bernard R. Ortiz de Montellano and Thelma Ortiz de Montellano.
Norman: University of Oklahoma Press, 1997.

\bibitem{taube1993}
Taube, Karl A.
\textit{Aztec and Maya Myths}.
Austin: University of Texas Press, 1993.

\bibitem{elliott2006}
Elliott, John H.
\textit{Empires of the Atlantic World: Britain and Spain in America, 1492--1830}.
New Haven: Yale University Press, 2006.

\bibitem{diamond1997}
Diamond, Jared.
\textit{Guns, Germs, and Steel: The Fates of Human Societies}.
New York: W. W. Norton, 1997.

\bibitem{scott1998}
Scott, James C.
\textit{Seeing Like a State: How Certain Schemes to Improve the Human Condition Have Failed}.
New Haven: Yale University Press, 1998.

\bibitem{taleb2012}
Taleb, Nassim Nicholas.
\textit{Antifragile: Things That Gain from Disorder}.
New York: Random House, 2012.

\bibitem{jacobson1995}
Jacobson, Ted.
``Thermodynamics of Spacetime: The Einstein Equation of State.''
\textit{Physical Review Letters} 75, no. 7 (1995): 1260--1263.

\bibitem{verlinde2011}
Verlinde, Erik.
``On the Origin of Gravity and the Laws of Newton.''
\textit{Journal of High Energy Physics} 2011, no. 4 (2011).

\bibitem{doctorow2023}
Doctorow, Cory.
\textit{The Internet Con: How to Seize the Means of Computation}.
London: Verso, 2023.

\end{thebibliography}

\end{document}
